\documentclass[11pt,a4paper]{article}

% ---------- Packages ----------
\usepackage[utf8]{inputenc}
\usepackage[T1]{fontenc}
\usepackage{lmodern}
\usepackage[margin=2.4cm]{geometry}
\usepackage{graphicx}
\usepackage{xcolor}
\usepackage{tabularx}
\usepackage{booktabs}
\usepackage{array}
\usepackage{longtable}
\usepackage{listings}
\usepackage{enumitem}
\usepackage{fancyhdr}
\usepackage{titlesec}
\usepackage[hidelinks]{hyperref}

% ---------- Colors ----------
\definecolor{techblue}{RGB}{0,51,102}
\definecolor{codebg}{RGB}{245,245,245}
\definecolor{codekw}{RGB}{0,90,160}

% ---------- Headers ----------
\pagestyle{fancy}
\fancyhf{}
\lhead{\textcolor{techblue}{ANN Accelerator Controller}}
\rhead{\textcolor{techblue}{User Manual}}
\cfoot{\thepage}
\renewcommand{\headrulewidth}{0.4pt}

\titleformat{\section}{\color{techblue}\Large\bfseries}{\thesection}{0.6em}{}
\titleformat{\subsection}{\color{techblue}\large\bfseries}{\thesubsection}{0.6em}{}

% ---------- Code style ----------
\lstdefinestyle{shell}{
  backgroundcolor=\color{codebg},
  basicstyle=\ttfamily\small,
  breaklines=true,
  frame=single,
  rulecolor=\color{gray!40},
  keywordstyle=\color{codekw}\bfseries,
  columns=fullflexible,
  showstringspaces=false
}
\lstdefinestyle{verilog}{
  backgroundcolor=\color{codebg},
  basicstyle=\ttfamily\footnotesize,
  breaklines=true,
  frame=single,
  rulecolor=\color{gray!40},
  keywordstyle=\color{codekw}\bfseries,
  morekeywords={module,input,output,logic,parameter,int,bit,string},
  showstringspaces=false
}

\newcolumntype{Y}{>{\raggedright\arraybackslash}X}

% ============================================================
\begin{document}

% ---------- Title Page ----------
\begin{titlepage}
\centering
\vspace*{1.5cm}
{\color{techblue}\Huge\bfseries ANN Accelerator Controller\par}
\vspace{0.5cm}
{\Large User Manual\par}
\vspace{0.3cm}
{\large A SystemVerilog Controller for a Y-Flash Memristor Array\par}
\vspace{2cm}
\rule{\textwidth}{0.8pt}\par
\vspace{0.6cm}
{\large
\begin{tabular}{rl}
\textbf{Team:} & Mohammad Abu Ammar \& Lana Bakrieh \\[4pt]
\textbf{Supervisor:} & Nir Ben-Haim \\[4pt]
\textbf{Institution:} & Technion --- VLSI Laboratory \\
\end{tabular}\par}
\vspace{0.6cm}
\rule{\textwidth}{0.8pt}\par
\vfill
{\large \today\par}
\end{titlepage}

\tableofcontents
\newpage

% ============================================================
\section{Project Goal \& Overview}

The \textbf{ANN Accelerator Controller} is a digital control unit, written in
SystemVerilog, that drives an Analog Neural Network (ANN) implemented on a
\textbf{Y-Flash memristor crossbar array}. It acts as the bridge between a host
processor and the analog in-memory-computing core.

The controller accepts five commands from a host through a parallel interface
and translates them into precisely timed \emph{pulse trains} that the analog
array requires:

\begin{itemize}[leftmargin=1.4em]
  \item \textbf{Programming} a quantized weight into a selected memristor cell.
  \item \textbf{Verifying} the programmed value with a closed feedback loop and
        automatically re-programming or erasing on mismatch.
  \item \textbf{Erasing} a cell.
  \item \textbf{Reading} a stored weight back from the array.
  \item \textbf{Inference}: collecting an $8\times8$ input image into an input
        buffer and streaming it bit-serially into the array for
        matrix--vector multiplication.
\end{itemize}

The central design challenge --- and the project's main contribution --- is the
\textbf{closed-loop program/verify algorithm}. Because memristor programming is
analog and imprecise, the controller reads back each weight after programming.
If the cell is \emph{under-programmed}, it re-applies program pulses (up to
\texttt{MAX\_PROG\_RETRIES}); if it is \emph{over-programmed}, it erases the
cell and retries. A lookup table (LUT) sets the number of program-pulse
repetitions per target weight value, allowing weight-dependent pulse shaping.

% ============================================================
\section{Architecture Components}

The system consists of three RTL modules and three supporting packages. The
top-level module is \texttt{ann\_controller}.

\begin{table}[h]
\centering
\renewcommand{\arraystretch}{1.3}
\begin{tabularx}{\textwidth}{l l Y}
\toprule
\textbf{Role} & \textbf{Module name} & \textbf{Source file} \\
\midrule
Parallel Interface & \texttt{parallel\_interface} &
  \texttt{rtl/Parallel\_Interface/parallel\_interface.sv} \\
ANN Controller (top) & \texttt{ann\_controller} &
  \texttt{rtl/Controller/controller.sv} \\
Input Buffer & \texttt{input\_buffer} &
  \texttt{rtl/Input\_Buffer/input\_buffer.sv} \\
\bottomrule
\end{tabularx}
\end{table}

\subsection{Parallel Interface (\texttt{parallel\_interface})}
Decodes the 32-bit host word \texttt{host\_data} into an 8-bit data byte and a
16-bit address, passes through the 3-bit command \texttt{host\_cmd}, and asserts
\texttt{valid} only when the command is active and the one-hot address tail is
well formed. The host word layout is
\texttt{host\_data[31:24]} = data byte and
\texttt{host\_data[23:0]} = one-hot \{PE, SA, column, row\}.

\subsection{ANN Controller (\texttt{ann\_controller})}
The heart of the system. It implements a 9-state main FSM
(\texttt{S\_IDLE, S\_RESET, S\_PROGRAM, S\_VERIFY, S\_ERASE, S\_READ,
S\_COLLECT\_DATA, S\_COMPUTE, S\_RESULT}) plus three sub-FSMs for the
\texttt{PROGRAM}, \texttt{VERIFY}, and \texttt{ERASE} sequences. It generates
the 3-bit \texttt{pulses} mode bus and the packed 32-bit \texttt{ann\_core\_word}
(data byte + one-hot address) toward the analog core, and orchestrates the
input buffer during inference.

\subsection{Input Buffer (\texttt{input\_buffer})}
A $64\times8$-bit memory. During inference it is loaded with an $8\times8$
image; during compute it outputs eight channels \texttt{D0--D7}, one bit per
channel per cycle (bit-serial, LSB-first). During weight verification it returns
the full byte at the addressed location.

% ============================================================
\section{Parameter Configuration}

All defaults are defined in \texttt{controller\_pkg} and the supporting
packages. The top module is parameterizable:

\begin{lstlisting}[style=verilog]
module ann_controller #(
    parameter int ADDR_WIDTH            = controller_pkg::DEFAULT_ADDR_WIDTH,   // 8
    parameter int WEIGHT_WIDTH          = controller_pkg::DEFAULT_WEIGHT_WIDTH, // 16
    parameter bit USE_WEIGHT_PULSE_LUT  = 1'b1,
    parameter string WEIGHT_PULSE_LUT_FILE =
        "target/Controller/programming_inputs/weight_pulse_lut.mem"
) ( ... );
\end{lstlisting}

\begin{longtable}{l l Y}
\toprule
\textbf{Parameter} & \textbf{Default} & \textbf{Description} \\
\midrule
\endhead
\texttt{ADDR\_WIDTH}            & 8   & Host address width (default). \\
\texttt{WEIGHT\_WIDTH}          & 16  & Weight data width (default). \\
\texttt{USE\_WEIGHT\_PULSE\_LUT}& 1   & Enable LUT-driven program-pulse repeat count. \\
\texttt{WEIGHT\_PULSE\_LUT\_FILE} & \emph{(path)} & \texttt{\$readmemh} LUT file: weight $\rightarrow$ repeat count. \\
\midrule
\multicolumn{3}{l}{\textit{Address map (\texttt{controller\_pkg}, total \texttt{WEIGHT\_ADDR\_WIDTH}=10)}} \\
\texttt{BLOCK\_ID\_WIDTH}     & 2 & PE / block select (\texttt{address[9:8]}). \\
\texttt{SUB\_BLOCK\_ID\_WIDTH}& 2 & Sub-array select (\texttt{address[7:6]}). \\
\texttt{COL\_ID\_WIDTH}       & 3 & Column within sub-block (\texttt{address[5:3]}). \\
\texttt{ROW\_ID\_WIDTH}       & 3 & Row within sub-block (\texttt{address[2:0]}). \\
\midrule
\multicolumn{3}{l}{\textit{Pulse generation (\texttt{controller\_pkg})}} \\
\texttt{TREAD}  & 2 & READ burst width (clock cycles). \\
\texttt{TPROG}  & 2 & PROGRAM burst width. \\
\texttt{TERASE} & 2 & ERASE burst width. \\
\texttt{TINF}   & 8 & INFERENCE burst width (min 8 for bit-serial). \\
\texttt{PULSE\_GAP} & 1 & Idle (HiZ) cycles between bursts. \\
\texttt{PULSE\_NUM\_*} & 1 & Number of bursts per operation. \\
\texttt{MAX\_PROG\_RETRIES} & 3 & Re-PROG attempts before falling back to ERASE. \\
\midrule
\multicolumn{3}{l}{\textit{Buffer / array sizing}} \\
\texttt{BUFFER\_SIZE}        & 64 & Input-buffer depth ($8\times8$ image). \\
\texttt{BUFFER\_DATA\_WIDTH} & 8  & Buffer word width. \\
\texttt{TOTAL\_WEIGHT\_LOCATIONS} & 1024 & $4\times4\times8\times8$ array cells. \\
\texttt{DEFAULT\_NUM\_WEIGHTS} & 640 & Default $10\times64$ weight matrix. \\
\bottomrule
\end{longtable}

% ============================================================
\section{Top-Level I/O Interface}

The following table lists every port of the top-level module
\texttt{ann\_controller} with its exact direction, bit-width, and function
inferred from the RTL.

\begin{longtable}{l c c Y}
\caption{\texttt{ann\_controller} port list}\\
\toprule
\textbf{Signal} & \textbf{Dir} & \textbf{Width} & \textbf{Description} \\
\midrule
\endfirsthead
\toprule
\textbf{Signal} & \textbf{Dir} & \textbf{Width} & \textbf{Description} \\
\midrule
\endhead

\multicolumn{4}{l}{\textit{Global}} \\
\texttt{clk}    & In & 1 & System clock. \\
\texttt{rst\_n} & In & 1 & Active-low asynchronous reset. \\
\midrule

\multicolumn{4}{l}{\textit{Controller $\leftrightarrow$ Parallel Interface}} \\
\texttt{valid}   & In & 1  & Asserted when a command/data beat is ready. \\
\texttt{data}    & In & 8  & Data byte from the parallel interface (weight/pixel payload). \\
\texttt{address} & In & 16 & Decoded address \{reserved[15:10], block[9:8], sub\_block[7:6], col[5:3], row[2:0]\}. \\
\texttt{cmd}     & In & 3  & Command (\texttt{CMD\_HIZ/READ/PROG/ERASE/INF}). \\
\midrule

\multicolumn{4}{l}{\textit{Controller $\leftrightarrow$ ANN Core}} \\
\texttt{ann\_reset}        & Out & 1  & Reset pulse toward the ANN core (asserted in \texttt{S\_RESET}). \\
\texttt{op\_done}          & In  & 1  & Core handshake: advances PROG/ERASE sub-FSMs. \\
\texttt{ann\_core\_word}   & Out & 32 & Packed bus: [31:24] data byte, [23:0] one-hot \{PE,SA,col,row\}. \\
\texttt{pulses}            & Out & 3  & Pulse-mode bus (HIZ/READ/PROG/ERASE/INF) gated by burst timing. \\
\texttt{weight\_read\_data}& In  & 4  & Quantized weight read back from the memristor during VERIFY. \\
\midrule

\multicolumn{4}{l}{\textit{Controller $\leftrightarrow$ Input Buffer}} \\
\texttt{buf\_reg\_add}    & Out & 6 & Buffer address. \\
\texttt{buf\_reg\_ctrl}   & Out & 3 & Buffer control (\texttt{CTRL\_IDLE/DATA\_LOAD/COMPUTE/RESULT\_OUT/WEIGHT\_READ}). \\
\texttt{buf\_read\_write} & Out & 1 & 1 = write to buffer, 0 = read. \\
\texttt{buf\_bit\_sel}    & Out & 3 & Bit index 0--7 for bit-serial D0--D7 output (LSB-first). \\
\texttt{buf\_data\_out}   & Out & 8 & Data written into the buffer (from host / parallel interface). \\
\texttt{buf\_ready}       & In  & 1 & Buffer ready/valid flag. \\
\texttt{buf\_data}        & In  & 8 & Full byte read from the buffer at the current address. \\
\midrule

\multicolumn{4}{l}{\textit{Status}} \\
\texttt{busy} & Out & 1 & High whenever the controller is not in \texttt{S\_IDLE}. \\
\bottomrule
\end{longtable}

\subsection{Command Encoding}
\begin{table}[h]
\centering
\renewcommand{\arraystretch}{1.2}
\begin{tabularx}{\textwidth}{l c Y}
\toprule
\textbf{Command} & \texttt{cmd[2:0]} & \textbf{Operation} \\
\midrule
\texttt{CMD\_HIZ}   & \texttt{000} & Idle / high-impedance, no pulses. \\
\texttt{CMD\_READ}  & \texttt{001} & Read a weight from the addressed cell. \\
\texttt{CMD\_PROG}  & \texttt{010} & Program a weight (LUT-shaped pulse train) then auto-verify. \\
\texttt{CMD\_ERASE} & \texttt{011} & Erase the addressed cell. \\
\texttt{CMD\_INF}   & \texttt{100} & Collect $8\times8$ data and run bit-serial inference. \\
\bottomrule
\end{tabularx}
\end{table}

% ============================================================
\section{Execution \& Simulation Guide}

The project is driven by \texttt{scripts/run\_sim.py}, which wraps ModelSim's
\texttt{vlog} (compile) and \texttt{vsim} (simulate). Run all commands from the
project root.

\subsection{Prerequisites}
\begin{itemize}[leftmargin=1.4em]
  \item ModelSim / QuestaSim ASE (Intel FPGA 17.0) with \texttt{vlog}/\texttt{vsim} on \texttt{PATH}.
  \item Python 3.8+ (plus \texttt{numpy}, \texttt{scikit-learn} for stimulus generation).
\end{itemize}

\subsection{List available modules and testbenches}
\begin{lstlisting}[style=shell]
python scripts/run_sim.py list
\end{lstlisting}

\subsection{Compile}
\begin{lstlisting}[style=shell]
# RTL only
python scripts/run_sim.py compile -m Controller -t rtl
# Full verification set (RTL + testbenches)
python scripts/run_sim.py compile -m Controller -t verif
\end{lstlisting}

\subsection{Run a single testbench (console / batch)}
\begin{lstlisting}[style=shell]
python scripts/run_sim.py sim -m Controller -tb controller_prog_verify_lut_tb
\end{lstlisting}

\subsection{Run with GUI + waveforms}
Passing \texttt{--do-file} opens the ModelSim GUI and loads the wave layout.
Use the \texttt{*\_waves\_tb} wrapper testbench with its matching \texttt{.do}:
\begin{lstlisting}[style=shell]
python scripts/run_sim.py sim -m Controller \
  -tb controller_prog_verify_lut_tb_waves_tb \
  --do-file tb/Controller/do/waves/controller_prog_verify_lut_tb_waves.do
\end{lstlisting}

\subsection{Run all testbenches for a module}
\begin{lstlisting}[style=shell]
python scripts/run_sim.py run_all -m Controller
python scripts/run_sim.py run_all -m Input_Buffer -m Parallel_Interface
\end{lstlisting}

\subsection{Clean generated artifacts}
\begin{lstlisting}[style=shell]
python scripts/run_sim.py clean -a              % work/ + target/ outputs
python scripts/run_sim.py clean -w              % work/ only
python scripts/run_sim.py clean -t -m Controller
\end{lstlisting}

\subsection{Detailed runbooks}
Per-testbench compile/run commands (batch and GUI wave) are maintained under
\texttt{docs/verification/}:
\begin{itemize}[leftmargin=1.4em]
  \item \texttt{TESTBENCH\_CATALOG.md} --- what each testbench checks
  \item \texttt{REGULAR\_TESTBENCH\_RUNBOOK.md} --- batch (\texttt{vsim -c}) runs
  \item \texttt{WAVE\_TESTBENCH\_RUNBOOK.md} --- GUI + \texttt{--do-file} wave runs
\end{itemize}

\subsection{Available testbenches}
\begin{longtable}{l Y}
\toprule
\textbf{Module} & \textbf{Testbenches} \\
\midrule
\endhead
Controller & \texttt{controller\_addr\_pulse\_tb},
             \texttt{controller\_prog\_verify\_lut\_tb},
             \texttt{controller\_prog\_verify\_lut\_10w\_tb},
             \texttt{parallel\_interface\_controller\_integration\_tb},
             \texttt{controller\_host\_erase\_tb},
             \texttt{controller\_inf\_buffer\_flow\_tb},
             \texttt{controller\_host\_read\_reorder\_tb} \\
\midrule
Input\_Buffer & \texttt{input\_buffer\_bit\_serial\_tb},
                \texttt{input\_buffer\_reset\_behavior\_tb},
                \texttt{input\_buffer\_full\_overwrite\_tb} \\
\midrule
Parallel\_Interface & \texttt{parallel\_interface\_extract\_tb} \\
\bottomrule
\end{longtable}

\vspace{1em}
\noindent\textbf{Output location.} Reports and logs are written under
\texttt{target/<module>/} (e.g.
\texttt{target/Controller/prog/prog\_verify\_report.txt}).

\end{document}
